% Spectre de la lumière solaire et quelques raies sombres de Fraunhofer ; au corrigé, élément responsable de chaque raie (exercice 15, chapitre 20)
\begin{tikzpicture}[x=0.3mm,y=1cm]
\shade[shading=otArcEnCiel] (400,0) rectangle (750,0.7);
\draw[otInk!50,line width=0.4pt] (400,0) rectangle (750,0.7);
\foreach \l in {656.3,589.0,527.0,517.3,486.1,434.0}{\draw[black,line width=1.1pt] (\l,0) -- (\l,0.7);}
\foreach \l/\e in {589.0/Na,517.3/Mg,434.0/H}{%
  \draw[otMuted,line width=0.4pt] (\l,0.7) -- (\l,0.85);
  \node[pcLegende,anchor=south,align=center] at (\l,0.85) {\num{\l}\solution{\\\textbf{\e}}};}
\foreach \l/\e in {656.3/H,527.0/Fe,486.1/H}{%
  \draw[otMuted,line width=0.4pt] (\l,0) -- (\l,-0.15);
  \node[pcLegende,anchor=north,align=center] at (\l,-0.15) {\num{\l}\solution{\\\textbf{\e}}};}
\node[pcLegende,anchor=west] at (755,0.35) {$\lambda$ (nm)};
\end{tikzpicture}
\par\smallskip\legende{\enonce{Spectre de la lumière du Soleil : longueurs d'onde (en nm) de six raies sombres.}\solution{Trois raies de l'hydrogène, une du sodium, une du fer et une du magnésium.}}
